\section*{第 \ref{ch:g12:proton-nmr} 章——质子核磁共振谱}

\begin{solution}{exo:g12:proton-nmr:1}
甲烷：1。乙烷：1（两个 \ce{CH3} 相同）。丙烷：2（两个
\ce{CH3} 相同，一个 \ce{CH2}）。甲醇：2（\ce{CH3} 和 \ce{OH}）。
甲氧基甲烷：1。
\end{solution}

\begin{solution}{exo:g12:proton-nmr:2}
2 个相邻\omterm{def:g9:inside-the-atom:proton}{质子}：三重峰，1:2:1。3 个相邻\omterm{def:g9:inside-the-atom:proton}{质子}：四重峰，1:3:3:1。
0 个相邻\omterm{def:g9:inside-the-atom:proton}{质子}：单峰。
\end{solution}

\begin{solution}{exo:g12:proton-nmr:3}
8 个\omterm{def:g9:inside-the-atom:proton}{质子}共计 \qty{40}{mm}：每个\omterm{def:g9:inside-the-atom:proton}{质子} \qty{5}{mm}。这些峰
分别代表 3、2 和 3 个\omterm{def:g9:inside-the-atom:proton}{质子}。
\end{solution}

\begin{solution}{exo:g12:proton-nmr:4}
醛基：9.7 至 \qty{10.0}{ppm}。链上的 \ce{CH3}：0.7 至
\qty{1.3}{ppm}。
\end{solution}

\begin{solution}{exo:g12:proton-nmr:5}
其 12 个\omterm{def:g12:proton-nmr:equivalent}{等效质子}给出一个尖锐、强度大的单峰，容易
将其定在零点。氘代\omterm{def:g6:solutions-and-solubility:solute}{溶剂}本身不给出\omterm{def:g9:inside-the-atom:proton}{质子}信号，
否则会掩盖样品（仅几毫克）的信号。
\end{solution}

\begin{solution}{exo:g12:proton-nmr:6}
两组信号。与氯相邻的 \ce{CH2}：2 个\omterm{def:g9:inside-the-atom:proton}{质子}，四重峰（3
个相邻\omterm{def:g9:inside-the-atom:proton}{质子}），位于 2.5--\qty{4.0}{ppm} 范围内。\ce{CH3}：3 个\omterm{def:g9:inside-the-atom:proton}{质子}，
三重峰（2 个相邻\omterm{def:g9:inside-the-atom:proton}{质子}），位于通常的 \ce{CH3} 范围附近，略高，
因为氯与之相隔两个\omterm{def:g7:atoms-and-molecules:atom}{原子}。
\end{solution}

\begin{solution}{exo:g12:proton-nmr:7}
三组信号：两个 \ce{CH3}，6 个\omterm{def:g12:proton-nmr:equivalent}{等效质子}，双峰（1
个相邻\omterm{def:g9:inside-the-atom:proton}{质子}，即 \ce{CH}）；\ce{CH}，1 个\omterm{def:g9:inside-the-atom:proton}{质子}，被 6 个相邻\omterm{def:g9:inside-the-atom:proton}{质子}裂分
为 7 重峰，位于 \ce{H-C-O} 范围（3.3--\qty{4.5}{ppm}）；
\ce{OH}：1 个\omterm{def:g9:inside-the-atom:proton}{质子}，单峰。
\end{solution}

\begin{solution}{exo:g12:proton-nmr:8}
(a) 丙酮；(b) 乙酸乙酯；(c) 丙酸。
\end{solution}

\begin{solution}{exo:g12:proton-nmr:9}
\ce{CH2} 与酯基的氧\omterm{def:g7:atoms-and-molecules:atom}{原子}相连，氧\omterm{def:g7:atoms-and-molecules:atom}{原子}\omterm{def:g11:polarity-and-cohesion:electronegativity}{电负性}大，
使其\omterm{def:g9:inside-the-atom:proton}{质子}去屏蔽，位于 \ce{H-C-O} 范围，
3.3--\qty{4.5}{ppm}。\ce{CH3} 离氧再远一个碳，仍处于
烷链范围内。
\end{solution}

\begin{solution}{exo:g12:proton-nmr:10}
丙酮：一个 \ce{C=O}（酮，\qty{1715}{cm^{-1}}）和六个\omterm{def:g12:proton-nmr:equivalent}{等效
质子}，给出一个单峰。丙醛会给出三组信号，其中一组
为 \ce{CH3-CH2-CHO} 的醛基\omterm{def:g9:inside-the-atom:proton}{质子}，位于 9.7--\qty{10.0}{ppm} 附近。
\end{solution}

\begin{solution}{exo:g12:proton-nmr:11}
\ce{OH} \omterm{def:g9:inside-the-atom:proton}{质子}在\omterm{def:g7:atoms-and-molecules:molecule}{分子}间快速交换，因此给出一个
单峰，且不使相邻\omterm{def:g9:inside-the-atom:proton}{质子}裂分。因此 \ce{CH2}
仅被 \ce{CH3} 的 3 个\omterm{def:g9:inside-the-atom:proton}{质子}裂分，为四重峰。
\end{solution}

\begin{solution}{exo:g12:proton-nmr:12}
丁酸 \ce{CH3-CH2-CH2-COOH}：4 组信号，三重峰（\ce{CH3}），一个
6 重峰（中间的 \ce{CH2}，5 个相邻\omterm{def:g9:inside-the-atom:proton}{质子}），三重峰
（\ce{CH2-C=O}），一个位于 11--\qty{12}{ppm} 附近的单峰。乙酸乙酯：一个
单峰、一个四重峰、一个三重峰。丙酸甲酯 \ce{CH3-CH2-CO-O-CH3}：
三重峰、四重峰（\ce{CH2-C=O}，2.0--\qty{2.4}{ppm}）、单峰
（\ce{O-CH3}）。甲酸丙酯 \ce{H-CO-O-CH2-CH2-CH3}：4 组信号，一个
位于最左侧的单峰（\ce{C=O} 碳上的 H），三重峰
（\ce{O-CH2}），6 重峰，三重峰（\ce{CH3}）。
\end{solution}

\begin{solution}{exo:g12:proton-nmr:13}
\ce{OH} 信号：\ce{O-H} 上的\omterm{def:g9:inside-the-atom:proton}{质子}与重水中的氘交换，
\ce{O-D} 不给出\omterm{def:g9:inside-the-atom:proton}{质子}信号。加入 \ce{D2O} 振荡后消失的信号，
属于连在 O 或 N 上的\omterm{def:g9:inside-the-atom:proton}{质子}。
\end{solution}

\begin{solution}{exo:g12:proton-nmr:14}
乙醇还给出一个位于 \qty{3.69}{ppm} 附近的四重峰，一个位于
\qty{1.23}{ppm} 附近的三重峰和一个 \ce{OH} 单峰。酯的单峰：\qty{30}{mm}
对应 3 个\omterm{def:g9:inside-the-atom:proton}{质子}，每个酯\omterm{def:g9:inside-the-atom:proton}{质子} \qty{10}{mm}。乙醇的三重峰：
\qty{3}{mm} 对应 3 个\omterm{def:g9:inside-the-atom:proton}{质子}，每个乙醇\omterm{def:g9:inside-the-atom:proton}{质子} \qty{1}{mm}。由于每个
\omterm{def:g7:atoms-and-molecules:molecule}{分子}在这些信号中各贡献一个 \ce{CH3}，两者的\omterm{def:g10:the-mole:amount}{物质的量}之比为
$1/10$：约每 10 个酯\omterm{def:g7:atoms-and-molecules:molecule}{分子}对应 1 个乙醇\omterm{def:g7:atoms-and-molecules:molecule}{分子}。
\end{solution}

\begin{solution}{exo:g12:proton-nmr:15}
四组信号：两个 \ce{CH3}（6 个\omterm{def:g9:inside-the-atom:proton}{质子}）为双峰；\ce{CH}（1
个\omterm{def:g9:inside-the-atom:proton}{质子}），被 6 + 2 = 8 个相邻\omterm{def:g9:inside-the-atom:proton}{质子}裂分为 9 重峰；\ce{CH2}（2
个\omterm{def:g9:inside-the-atom:proton}{质子}），与 \ce{O} 相邻，为双峰，位于 \ce{H-C-O} 范围；
\ce{OH}（1 个\omterm{def:g9:inside-the-atom:proton}{质子}）为单峰。\ce{CH} 信号的峰数最多。
\end{solution}

\begin{solution}{pb:g12:proton-nmr:1}
\textbf{1.} 含四个碳的\omterm{def:g11:organic-skeletons:alkane}{烷烃}为 \ce{C4H10}；一个 \ce{C=O}
减少两个氢，两个氧\omterm{def:g7:atoms-and-molecules:atom}{原子}不增加氢：\ce{C4H8O2}，
即四个碳的酸和酯的\omterm{def:g11:organic-skeletons:formulas}{分子式}。

\textbf{2.} 丁酸 \ce{CH3-CH2-CH2-COOH}；乙酸乙酯
\ce{CH3-CO-O-CH2-CH3}；丙酸甲酯 \ce{CH3-CH2-CO-O-CH3}；
甲酸丙酯 \ce{H-CO-O-CH2-CH2-CH3}。

\textbf{3.} 三种酯都含有酯基；这四种物质都含有一个
\ce{C=O}。

\textbf{4.} \ce{C=O}，位于酯的范围内（约 \qty{1735}{cm^{-1}}）。

\textbf{5.} 在 2500 至
\qty{3300}{cm^{-1}} 之间没有很宽的 \ce{O-H} 吸收带，因此不是丁酸。

\textbf{6.} 不能：三种酯含有相同的基团。

\textbf{7.} 三组。

\textbf{8.} 乙基 \ce{CH3-CH2-}：\ce{CH2} 被 3 个\omterm{def:g9:inside-the-atom:proton}{质子}裂分（四重峰），
\ce{CH3} 被 2 个\omterm{def:g9:inside-the-atom:proton}{质子}裂分（三重峰）。

\textbf{9.} 其\omterm{def:g9:inside-the-atom:proton}{质子}的相邻\omterm{def:g7:atoms-and-molecules:atom}{原子}上没有\omterm{def:g9:inside-the-atom:proton}{质子}：即
一个 \ce{CH3} 与 \ce{C=O} 碳或酯氧相连。

\textbf{10.} \ce{CH2} 与氧相连（\ce{H-C-O} 范围）。

\textbf{11.} 一个三重峰（\ce{CH3}），一个位于 2.0--\qty{2.4}{ppm} 附近的四重峰
（\ce{CH2} 紧邻 \ce{C=O}）以及一个位于 \ce{H-C-O} 范围的单峰
（\ce{O-CH3}）。四重峰将远低于 \qty{4.12}{ppm}，而单峰
远高于 \qty{2.04}{ppm}：因此不是未知物。

\textbf{12.} 共四组信号（最左侧的单峰、三重峰、6 重峰、
三重峰）：不是未知物，未知物只有三组信号。

\textbf{13.} 乙酸乙酯，\ce{CH3-CO-O-CH2-CH3}：单峰
\qty{2.04}{ppm} = \ce{CH3-C=O}；四重峰 \qty{4.12}{ppm} = \ce{O-CH2}；
三重峰 \qty{1.26}{ppm} = 末端 \ce{CH3}。

\textbf{14.} 单峰 3，四重峰 2，三重峰 3。

\textbf{15.} 每个\omterm{def:g9:inside-the-atom:proton}{质子} $72 / 8 = \qty{9}{mm}$。

\textbf{16.} 单峰 $3 \times 9 = \qty{27}{mm}$；三重峰
\qty{27}{mm}。

\textbf{17.} $18 + 27 + 27 = \qty{72}{mm}$。

\textbf{18.} 可以用嗅觉辨别：小\omterm{def:g7:atoms-and-molecules:molecule}{分子}酯以其果香著称。

\textbf{19.} 四重峰的台阶为 \qty{72}{mm} 中的 \qty{18}{mm}。
\end{solution}
